% The look shared by the Axon User Manual (axon-manual.tex) and the Axon
% Tutorial (../tutorial/axon-tutorial.tex): packages, page layout, the
% colours of the IDE's light theme, the Axon listings language and the
% environments and commands for code. Each document \input's it after its
% \documentclass.

\usepackage[margin=2.5cm]{geometry}
\usepackage[T1]{fontenc}
\usepackage[utf8]{inputenc}
\usepackage{lmodern}
\usepackage[scaled=0.88]{beramono}
\usepackage{xcolor}
\usepackage{listings}
\usepackage{graphicx}
\usepackage{booktabs}
\usepackage{longtable}
\usepackage{tabularx}
\usepackage{array}
\usepackage{amsmath}
\usepackage[shortlabels]{enumitem}
\usepackage{caption}
\usepackage[colorlinks=true,linkcolor=blue!50!black,urlcolor=blue!50!black,
            citecolor=blue!50!black]{hyperref}

\setlist{itemsep=2pt,topsep=4pt}
\setlength{\parskip}{4pt}
\setlength{\parindent}{0pt}
\setlength{\emergencystretch}{3em}
\captionsetup{font=small,labelfont=bf}

% Colours of the Axon IDE's light theme.
\definecolor{axkeyword}{HTML}{0033B3}
\definecolor{axbuiltin}{HTML}{00796B}
\definecolor{axstring}{HTML}{A31515}
\definecolor{axnumber}{HTML}{098658}
\definecolor{axcomment}{HTML}{6A737D}
\definecolor{axdoc}{HTML}{8A5A00}
\definecolor{axback}{HTML}{F6F8FA}
\definecolor{axframe}{HTML}{D0D7DE}
\definecolor{outback}{HTML}{FFFFFF}

\lstdefinelanguage{Axon}{
  morekeywords={function,end,return,if,elseif,else,for,in,while,break,
    continue,local,global,try,catch,import,from,as,export,extern,record,
    true,false,nothing,and,or,not},
  morekeywords=[2]{zeros,ones,eye,rand,linspace,size,sum,prod,mean,all,any,
    dot,adjoint,isequal,keep,eltype,reshape,sparse,speye,nnz,full,splu,
    append,get,keys,haskey,dict,print,println,error,string,length,typeof,
    sqrt,exp,log,log10,log2,sin,cos,tan,asin,acos,atan,atan2,deg2rad,rad2deg,
    sinh,cosh,tanh,asinh,acosh,atanh,
    abs,int,floor,ceil,round,real,imag,conj,arg,complex,div,mod,min,max,
    factorial,format,pi,sort,sortperm,atexit,range,isa,tr,trace,split,join,
    trim,ltrim,rtrim,upper,lower,contains,startswith,endswith,find,replace,
    repeat,parse_function,isapprox,input,vector},
  sensitive=true,
  morecomment=[l]{\#},
  morestring=[b]",
}

\lstdefinestyle{axonstyle}{
  language=Axon,
  basicstyle=\ttfamily\small,
  keywordstyle=\color{axkeyword}\bfseries,
  keywordstyle=[2]\color{axbuiltin},
  stringstyle=\color{axstring},
  commentstyle=\color{axcomment}\itshape,
  backgroundcolor=\color{axback},
  frame=single, rulecolor=\color{axframe}, framesep=4pt,
  xleftmargin=6pt, xrightmargin=6pt,
  columns=fullflexible, keepspaces=true,
  showstringspaces=false, upquote=true,
  aboveskip=8pt, belowskip=6pt,
  breaklines=true, breakatwhitespace=true,
  tabsize=2,
  literate={*}{{\char42}}1 {-}{{\char45}}1,
}
\lstdefinestyle{outstyle}{
  language={}, keywordstyle={}, stringstyle={}, commentstyle={},
  literate={*}{{\char42}}1 {-}{{\char45}}1,
  basicstyle=\ttfamily\small\color{black!75},
  backgroundcolor=\color{outback},
  frame=leftline, rulecolor=\color{axframe}, framerule=2pt,
  xleftmargin=10pt, framesep=6pt,
  columns=fullflexible, keepspaces=true, upquote=true,
  aboveskip=0pt, belowskip=8pt,
  escapeinside={(*@}{@*)},
}
\lstdefinestyle{shellstyle}{
  language={}, keywordstyle={}, stringstyle={}, commentstyle={},
  basicstyle=\ttfamily\small,
  backgroundcolor=\color{axback},
  frame=single, rulecolor=\color{axframe}, framesep=4pt,
  xleftmargin=6pt, xrightmargin=6pt,
  columns=fullflexible, keepspaces=true, upquote=true,
}
% axon: a complete program (checked by running it). axonout: what it prints.
% axonerr: a program that stops with an error; the following axonout is the
% error report. shell: a command line.
\lstnewenvironment{axon}{\lstset{style=axonstyle}}{}
\lstnewenvironment{axonerr}{\lstset{style=axonstyle}}{}
\lstnewenvironment{axonout}{\lstset{style=outstyle}}{}
\lstnewenvironment{shell}{\lstset{style=shellstyle}}{}
% axonrepl: a REPL session, checked with axon --transcript.
\lstnewenvironment{axonrepl}{\lstset{style=outstyle,aboveskip=8pt}}{}
\lstset{style=axonstyle}
\newcommand{\ex}[1]{\lstinputlisting[style=axonstyle,
  title={\small\ttfamily\detokenize{#1}}]{examples/#1}}
% Inline code, read verbatim: \code{x + 1}, or \code|{a, b}| with braces inside.
\newcommand{\code}{\lstinline[basicstyle=\ttfamily,keywordstyle=,
  keywordstyle={[2]},stringstyle=,commentstyle=,
  literate={*}{{\char42}}1 {-}{{\char45}}1]}
\newcommand{\menu}[1]{\textsf{#1}}
\newcommand{\key}[1]{\fbox{\footnotesize\sffamily #1}}
\newcommand{\fig}[3]{\begin{figure}[htbp]\centering
  \includegraphics[width=#2\linewidth]{figures/#1}\caption{#3}\label{fig:#1}
  \end{figure}}
